\documentclass[11pt,a4paper]{article}
\usepackage[T1]{fontenc}
\usepackage[utf8]{inputenc}
\usepackage{lmodern}
\usepackage[a4paper,margin=2cm]{geometry}
\usepackage{amsmath}
\usepackage{booktabs,tabularx}
\usepackage{enumitem}
\usepackage[dvipsnames]{xcolor}
\setlist{nosep,leftmargin=*}
\setlength{\parindent}{0pt}
\setlength{\parskip}{0.4em}
\pagestyle{plain}

\newcommand{\todo}[1]{\textcolor{BrickRed}{[#1]}}

\begin{document}

% ===================== PAGE 1: INVESTMENT POLICY (1 page) =====================
\begin{center}
{\Large\bfseries \ Matterhorn Capital Management\\[0.2em]
{\small AMC 2026, Round 1 \textbullet{} \ Group 3: Arauz, Di Pietro, Musi, Rocco, Rotondò, Selmi}
\end{center}

\section*{Investment policy}

\subsection*{1. Universe}
Long-only ETFs on broad asset classes; no leverage; few positions; cash earns no interest.
\todo{List of the ETFs the allocation may hold, with one-word role each. i.e. VT (global equity), GOVT (Treasuries), STIP (inflation protection), GLD (gold), BCI (commodities), BIL (cash substitute).}

\subsection*{2. Objective and risk budget}
\begin{itemize}
  \item Objective: maximize the Sharpe Ratio with an allocation that performs acceptably
      in every macro regime 
  \item Volatility target: 8\% p.a.\ (range 7--9\%)
  \item Ex ante volatility of our allocation: \textbf{\todo{x.x}\%}
  \item One-month 5\% VaR: \textbf{\todo{x.x}\% = USD \todo{x.x} million}
  \item Maximum acceptable drawdown: \textbf{\todo{$-$xx}\%}
\end{itemize}

\subsection*{3. Constraints, with reasons}
\begin{tabularx}{\textwidth}{@{}p{0.32\textwidth}X@{}}
\toprule
Constraint & Reason (today's model errors) \\
\midrule
Concentration: \todo{e.g.\ max xx\% per ETF} & \todo{reason} \\
Regions: \todo{e.g.\ \ldots}                 & \todo{reason} \\
Liquidity: \todo{e.g.\ \ldots}               & \todo{reason} \\
\bottomrule
\end{tabularx}

\subsection*{4. Risk contributions}
\begin{center}
\begin{tabular}{@{}lrr@{}}
\toprule
ETF & Capital weight & Risk contribution \\
\midrule
\todo{Ticker} & \todo{xx.x}\% & \todo{xx.x}\% \\
\todo{Ticker} & \todo{xx.x}\% & \todo{xx.x}\% \\
\todo{Ticker} & \todo{xx.x}\% & \todo{xx.x}\% \\
\todo{Ticker} & \todo{xx.x}\% & \todo{xx.x}\% \\
\todo{Ticker} & \todo{xx.x}\% & \todo{xx.x}\% \\
\todo{Ticker} & \todo{xx.x}\% & \todo{xx.x}\% \\
\todo{Ticker} & \todo{xx.x}\% & \todo{xx.x}\% \\
Cash          & \todo{x.x}\%  & 0.0\%         \\
\midrule
Total & 100.0\% & 100.0\% \\
\bottomrule
\end{tabular}
\end{center}
Position above half of the risk: \todo{none / ETF -- intended because \ldots}

\newpage
% ============ COMPARISON WITH THE ROBUST APPROACHES (1/2 page) ============
\section*{Our allocation vs.\ the robust approaches}

Equal weight: $w_i = 1/N$. Risk parity: $w_i \propto 1/\sigma_i$, volatilities from our own data (\todo{period}).

\begin{center}
\small
\begin{tabular}{@{}lrrrrrr@{}}
\toprule
 & \multicolumn{2}{c}{Our SAA} & \multicolumn{2}{c}{Equal weight} & \multicolumn{2}{c}{Risk parity} \\
\cmidrule(lr){2-3}\cmidrule(lr){4-5}\cmidrule(lr){6-7}
ETF & Weight & Risk contr. & Weight & Risk contr. & Weight & Risk contr. \\
\midrule
\todo{Ticker} & & & & & & \\
\todo{Ticker} & & & & & & \\
\todo{Ticker} & & & & & & \\
\todo{Ticker} & & & & & & \\
\todo{Ticker} & & & & & & \\
\todo{Ticker} & & & & & & \\
\todo{Ticker} & & & & & & \\
\midrule
Volatility & \multicolumn{2}{c}{\todo{x.x}\%} & \multicolumn{2}{c}{\todo{x.x}\%} & \multicolumn{2}{c}{\todo{x.x}\%} \\
\bottomrule
\end{tabular}
\end{center}

\todo{One paragraph: where the rules disagree with our weights, and what justifies our weights.}

\vspace{1.5em}
% ======================= BACKTEST PROTOCOL =======================
\section*{Backtest protocol}

\subsection*{1. Design}
\begin{itemize}
  \item Universe: \todo{ETFs (and proxies used for the missing months)}.
  \item Rule that determines the weights: \todo{method; estimation window, e.g.\ rolling 60 months}.
  \item Rebalancing: monthly, to the SAA weights.
  \item Sample period: \todo{start -- end}; data: \todo{Bloomberg, monthly total returns in USD}.
\end{itemize}

\subsection*{2. Validation}
\begin{itemize}
  \item Data used for tuning: \todo{start -- end}.
  \item Data held out for validation: \todo{start -- end}, evaluated once, after the choice.
  \item Number of variations tried: \textbf{\todo{N}} \todo{(methods, universes, constraints, parameters)}.
\end{itemize}

\subsection*{3. Risk}
\begin{center}
\begin{tabular}{@{}lrr@{}}
\toprule
 & Tuning sample & Hold-out sample \\
\midrule
Mean return, annualised (t-statistic) & \todo{x.x}\% (\todo{x.x}) & \todo{x.x}\% (\todo{x.x}) \\
Volatility, annualised                & \todo{x.x}\% & \todo{x.x}\% \\
Maximum drawdown                      & \todo{$-$xx.x}\% & \todo{$-$xx.x}\% \\
One-month 5\% VaR                     & \todo{x.x}\% & \todo{x.x}\% \\
Sharpe Ratio                          & \todo{x.xx} & \todo{x.xx} \\
\bottomrule
\end{tabular}
\end{center}

\subsection*{4. The discount}
Backtested Sharpe Ratio: \textbf{\todo{x.xx}}. Discounted value we consider realistic: \textbf{\todo{x.xx}}.\\
Justification: \todo{rule of thumb of Class 4 (divide by 2), larger discount for more tunable parameters / variations tried;
comparison with the hold-out Sharpe Ratio.}

\end{document}
